\scanpage{26}
\mspage{6}
notons que pour des $X$ lisses et propres, on a bien le formalisme de dualité habituel (autodualité de $H_{\mathrm{DR}}(X)$, \inserted{d'où} \unsure{hom.} de Gysin, classe de coh. associée à un cycle alg, \deleted{\unsure{et}} compatibilité avec images directes, inverses, intersections) $\ldots$ le tout se faisant de façon tout analogue à ce que, dans le point de vue ``cohomologie de Hodge'' j'avais fait au Sém. Bourbaki.
Mais le ``lecteur informé'' au courant des histoires $f_!,f^!$ verra à quel point une \inserted{simple} dualité de ce type est \unsure{candide} et primaire.
On aimerait un formalisme $f_!,f^!$ du type de celui développé en long et en large dans le livre Hartshorne (Springer), \deleted{\unsure{et}} pour des morphismes $\pm$ quelconques, \inserted{disons \deleted{\unsure{propres,}}\inserted{séparés de} type fini;} de schémas.
Un tel formalisme, \inserted{impliquant des coeff. tels les complexes de D.R.,} devrait alors automatiquement donner \unsure{égal$^{\mathrm t}$} une description de l'\uline{homologie} singulière (comme il est suggéré par ailleurs dans mes papiers bleus$\ldots$).
Il y aurait lieu en particulier d'étendre ce formalisme, du cas de complexes à différentielles linéaires, au cas des complexes
\margin{25}
